\documentclass[11pt]{article}
\usepackage[margin=1in]{geometry}
\usepackage[T1]{fontenc}
\usepackage{lmodern,amsmath,amssymb,amsthm,microtype}
\usepackage[colorlinks=true,linkcolor=blue,citecolor=blue,urlcolor=blue]{hyperref}
\hypersetup{pdftitle={A simple brace outside asymmetric products of two trivial braces},
pdfauthor={Alec Kriebel},pdfsubject={A negative answer to Cedo's Problem 9},
pdfkeywords={left brace, simple brace, asymmetric product, metabelian group}}
\newtheorem{theorem}{Theorem}
\newtheorem{lemma}[theorem]{Lemma}
\newtheorem{proposition}[theorem]{Proposition}
\newcommand{\F}{\mathbb F}
\newcommand{\C}{\mathrm C}
\newcommand{\id}{\operatorname{id}}
\newcommand{\wt}{\operatorname{wt}}
\newcommand{\Syl}{\operatorname{Syl}}
\title{A simple brace outside asymmetric products\\of two trivial braces}
\author{Alec Kriebel\\
\small\href{https://orcid.org/0009-0001-9320-500X}{ORCID: 0009-0001-9320-500X}}
\date{17 October 2026 \quad Version 1.0.0}
\begin{document}
\maketitle
\begin{abstract}
We give a negative answer to Ced\'o's Problem~9 in Oberwolfach Report
51/2019. An explicit left brace of order $20\,000$ is nontrivial and
simple, has metabelian multiplicative group with abelian Sylow
subgroups, and is not an asymmetric product of two trivial braces.
The obstruction is intrinsic: every such two-factor product has zero
triple star products within each additive Sylow subbrace, whereas the
constructed brace has $(x*x)*x=x^3\ne0$ in its $5$-primary subbrace.
The obstruction covers the full admissible symmetric-cocycle definition,
not only a chosen bilinear construction family. All operations and the
simplicity argument are given directly; exact verification programs and
independent adversarial audits accompany the note.
\end{abstract}

\section{The question}
A left brace $(B,+,\circ)$ has an abelian additive group, a multiplicative
group, and the compatibility
\[
 a\circ(b+c)=a\circ b-a+a\circ c.
\]
Write $\lambda_a(b)=-a+a\circ b$ and $a*b=\lambda_a(b)-b$.
The maps $\lambda_a$ are additive automorphisms and satisfy
$\lambda_{a\circ b}=\lambda_a\lambda_b$.
The brace is trivial when $a*b=0$ for all $a,b$.
An ideal is an additive subgroup invariant under all $\lambda_a$ and
normal in the multiplicative group; a nonzero brace is simple if its
only ideals are $0$ and itself.

Ced\'o asks whether every finite nontrivial simple left brace with a
metabelian multiplicative group and abelian multiplicative Sylow
subgroups is an asymmetric product of two trivial braces
\cite[p.~3228, Problem~9]{OWR}. We retain all these hypotheses.

\begin{theorem}\label{thm:main}
There is a nontrivial simple left brace $B$ of order $20\,000$ with
\[
 (B,+)\cong\C_5^4\times\C_2^5,\qquad
 (B,\circ)\cong(\C_5^3\rtimes_\tau\C_2)
                 \times(\C_2^4\rtimes_M\C_5),
\]
where $\tau$ negates the first and third coordinates of $\F_5^3$ and
$M$ is the five-cycle on the even-weight subspace of $\F_2^5$.
Its multiplicative group is metabelian with every Sylow subgroup
abelian. No brace isomorphic to $B$ is an asymmetric product of two
trivial braces.
\end{theorem}

\section{An invariant of every two-trivial-factor product}
We use the full asymmetric-product definition of
Catino--Colazzo--Stefanelli, in its left-brace formulation
\cite{CCS,BCJO,CJOembed}. With trivial factors $T,S$, an admissible
presentation has operations
\begin{align}
 (t,s)+(u,v)&=(t+u,s+v+b(t,u)),\label{eq:asymadd}\\
 (t,s)\circ(u,v)&=(t+\alpha_s(u),s+v),\label{eq:asymmult}
\end{align}
where $\alpha:S\to\operatorname{Aut}(T,+)$ is a group homomorphism.
Initially $b:T\times T\to S$ is only a normalized symmetric $2$-cocycle:
$b(t,0)=0$, $b(t,u)=b(u,t)$, and
\[
 b(t+u,w)+b(t,u)=b(t,u+w)+b(u,w).
\]
Admissibility means that these operations define a left brace, as in
the cited general definition.

\begin{lemma}\label{lem:obstruction}
If a finite brace has a presentation \eqref{eq:asymadd}--\eqref{eq:asymmult}
with $T,S$ trivial and all its multiplicative Sylow subgroups abelian,
then
\[
                     (a*b)*c=0
\]
for all $a,b,c$ in any one additive Sylow subbrace.
\end{lemma}
\begin{proof}
The cocycle identity gives
\[
 (t,s)-(u,v)=(t-u,s-v-b(t-u,u)),
\]
and therefore, before assuming biadditivity,
\[
 \lambda_{(t,s)}(u,v)=(\alpha_s(u),v-b(\alpha_s(u),t)).
\]
Since $(t,0)\circ(u,0)=(t+u,0)$, lambda composition evaluated on
$(w,0)$ forces $b(w,t+u)=b(w,t)+b(w,u)$. Symmetry makes $b$ biadditive.

For a prime $p$, write $T_p,S_p$ for the primary subgroups of the
finite abelian factor groups. Mixed-primary values of $b$ vanish, and
$b(T_p,T_p)\subseteq S_p$: each value is killed by the relevant prime
powers. Consequently the set $P=T_p\times S_p$ is an additive subgroup
of the full $p$-part of the order. Its additive law need not split,
especially for $p=2$; its order suffices. The action preserves $T_p$,
so the same set is also a multiplicative Sylow subgroup.

Commuting $(t,0)$ and $(0,s)$ inside this abelian multiplicative
Sylow subgroup gives $\alpha_s(t)=t$ for $s\in S_p,t\in T_p$.
For $a=(t,s)$ and $b_0=(u,v)$ in $P$, twisted additive subtraction
now gives
\[
 a*b_0=(u,v-b(u,t))-(u,v)=(0,-b(u,t)).
\]
Every $(0,w)$ with $w\in S_p$ has identity lambda action on $P$.
Thus $(a*b_0)*c=0$ for every $c\in P$.
\end{proof}

The lemma does not require simplicity or metabelianity.
Brace isomorphisms preserve additive Sylow subgroups and both star
operations, so a nonzero primary triple star excludes every possible
presentation, including one with different factors or coordinates.

\section{The explicit brace}
Let $R=x\F_5[x]/(x^5)$, with basis $x,x^2,x^3,x^4$.
In the ambient unital truncated polynomial ring, define
\[
 L(r)=r-\frac{r^2}{2}+\frac{r^3}{3}-\frac{r^4}{4},
 \qquad E(t)=t+\frac{t^2}{2}+\frac{t^3}{6}+\frac{t^4}{24}.
\]
All denominators are invertible in $\F_5$. Formal expansion through
degree four, with higher total degrees zero, proves
\begin{equation}\label{eq:log}
 L(E(t))=t,\quad E(L(r))=r,\quad
 L(r+u+ru)=L(r)+L(u).
\end{equation}
There is no division by five or infinite-series assumption.
Let $\sigma$ be substitution $x\mapsto-x$, and put
\[
                           \eta(r)=[x^4]L(r).
\]
Then $\eta(r+u+ru)=\eta(r)+\eta(u)$,
$\eta(\sigma r)=\eta(r)$, and $L(\sigma r)=\sigma L(r)$.

Let $V$ be the even-weight subspace of $\F_2^5$. For $v,w\in V$ set
\[
 \beta(v,w)=\sum_{i=0}^4v_iw_i,\qquad
 Q(v)=\frac{\wt(v)}2\pmod2.
\]
The integer weight identity proves
\begin{equation}\label{eq:quadratic}
 Q(v+w)=Q(v)+Q(w)+\beta(v,w).
\end{equation}
The restricted dot form $\beta$ is nondegenerate: the orthogonal
complement of the even-weight space in $\F_2^5$ consists of constant
vectors, and its only even-weight constant vector is zero.
Cyclic rotation $M$ preserves $V,\beta,Q$ and has order five.
Its fixed vectors in $V$ are zero, so $M-\id$ is invertible.

Take $B=R\times V\times\F_2$ with componentwise addition. For
$a=(r,v,z)$ write $s(a)=\eta(r)$ and $e(a)=z+Q(v)$, and define
\begin{equation}\label{eq:direct}
 (r,v,z)\circ(u,w,Z)=
 \bigl(r+\sigma^{e(a)}u+r\sigma^{e(a)}u,\
       v+M^{s(a)}w,\
       z+Z+\beta(v,M^{s(a)}w)\bigr).
\end{equation}
Exponents of $M$ are interpreted modulo five.

\begin{proposition}\label{prop:brace}
These operations make $B$ a left brace with the two groups stated
in Theorem~\ref{thm:main}.
\end{proposition}
\begin{proof}
Write $L(r)=l_1x+l_2x^2+l_3x^3+s x^4$ and $k=(l_1,l_2,l_3)$.
The map
\[
 \Psi(r,v,z)=((k,z+Q(v)),(v,s))
\]
is a bijection onto
$(\F_5^3\rtimes_\tau\C_2)\times(V\rtimes_M\C_5)$:
recover $r$ with $E$, and then $z$.
For a second input with corresponding coordinates $h,f,w,t$,
\eqref{eq:log} gives the output logarithm $L(r)+\sigma^{e(a)}L(u)$.
Its first three coefficients are $k+\tau^{e(a)}h$ and its fourth is
$s+t$. Equation~\eqref{eq:quadratic} gives
\[
 z+Z+\beta(v,M^s w)+Q(v+M^s w)
       =z+Q(v)+Z+Q(w)=e(a)+f.
\]
Thus the complete transformed product is exactly
$((k+\tau^{e(a)}h,e(a)+f),(v+M^s w,s+t))$.
This proves the group law, including associativity and inverses.

Direct additive subtraction in \eqref{eq:direct} yields
\begin{equation}\label{eq:lambda}
 \lambda_{(r,v,z)}(u,w,Z)=
 \bigl((1+r)\sigma^{z+Q(v)}u,\
       M^{\eta(r)}w,\
       Z+\beta(v,M^{\eta(r)}w)\bigr).
\end{equation}
The first block is an invertible $\F_5$-linear map, since $1+r$ is
a unit. The last two coordinates form an invertible triangular
$\F_2$-linear map. Every lambda map is therefore an additive
automorphism, proving the left-brace identity.
Finally $|B|=5^4\,2^4\,2=20\,000$.
\end{proof}

The commutator subgroup of the displayed multiplicative group is
\[
 (\tau-\id)\F_5^3\times(M-\id)V
               \cong\C_5^2\times\C_2^4,
\]
which is abelian of order $400$. Both actions are nontrivial, so the
group is nonabelian and metabelian. Explicit Sylow subgroups are
$\F_5^3\times\C_5$ and $\C_2\times V$ of orders $625$ and $32$,
respectively. Both are abelian, and Sylow conjugacy covers all Sylow
subgroups. No normal-Sylow assumption is used.

\section{Simplicity and exclusion}
\begin{proposition}\label{prop:simple}
Every nonzero ideal of $B$ is $B$.
\end{proposition}
\begin{proof}
For an ideal $I$, $i\in I$ and $g\in B$, both
$\lambda_g(i)-i$ and $i*g$ lie in $I$. The second inference follows
in the quotient brace $B/I$, where $i$ becomes zero; it is not an
assumption about arbitrary left ideals.

If a nonzero element $(r,v,z)\in I$ has $r\ne0$, its additive
sixfold multiple is $(r,0,0)\ne0$. Let $c x^j$ be the lowest nonzero
term of $r$. For $j<4$, formula \eqref{eq:lambda} gives
\[
 (\lambda_{(x^{4-j},0,0)}-\id)(r,0,0)=(c x^4,0,0).
\]
For $j=4$ this already holds without applying a lambda difference.
Additive rescaling gives $a=(x^4,0,0)\in I$.

Otherwise $I$ has a nonzero element $(0,v,z)$. If $v=0$ it is
$b=(0,0,1)$. If $v\ne0$, choose $w$ with $\beta(w,v)=1$.
Then $(\lambda_{(0,w,0)}-\id)(0,v,z)=b$.
Thus every nonzero ideal contains $a$ or $b$.

From $b$ one obtains
$b*(x,0,0)=(3x,0,0)$ and hence a nonzero ring element, which gives $a$.
Conversely, $\eta(x^4)=1$ implies
\[
 a*(0,w,0)=(0,(M-\id)w,0).
\]
Surjectivity of $M-\id$ puts every pure $V$ vector into $I$, and a
binary lambda difference gives $b$. Thus either starting point gives
all of $P_2=\{0\}\times V\times\F_2$ and also $(x,0,0)$.
For $j=1,2,3$,
\[
 (\lambda_{(x,0,0)}-\id)(x^j,0,0)=(x^{j+1},0,0).
\]
These monomials generate $P_5=R\times\{0\}\times\{0\}$ additively.
Hence $I$ contains both primary parts and equals $B$.
\end{proof}

On $P_5$, star multiplication is ordinary ring multiplication.
In particular $x*x=x^2\ne0$, so the brace is nontrivial, and
\[
                        (x*x)*x=x^3\ne0.
\]
Lemma~\ref{lem:obstruction} rules out every asymmetric product of two
trivial braces. Together with Propositions~\ref{prop:brace}
and~\ref{prop:simple}, this proves Theorem~\ref{thm:main} and gives the
full negative answer to the original question.

\section{Prior work, verification, and review status}
Asymmetric and matched products, radical-ring braces, truncated
logarithm/exponential identities and cyclic orthogonal actions are
established methods \cite{CCS,BCJO,CJOabundance,CJOembed}.
General matched-product formulas \cite{Bachiller} recover \eqref{eq:direct}, but their
existence statement does not replace the all-ideal simplicity proof.
No new general construction formalism is claimed. Recent work also
uses local commutative algebras and invariant homomorphisms
\cite{Verbeken}; those ingredients themselves are not a novelty claim.

A literature audit on 17 October 2026 found no prior full answer or
verified published full implication for this target. This is a bounded
finding, not historical-priority certification. It included the current
arXiv version of \cite{CO} and the recent constructions
\cite{DameleKT,DameleIB}. The full bodies of \cite{Rump2020,Rump2022}
and several final journal versions, including that of \cite{CO}, were
unavailable; their broader consequences are not excluded by abstracts.
The companion audit records exact versions, inspection boundaries,
access failures and citation-coverage limits.

The source/verification archive contains the original full proof,
independently authored adversarial reports, four exact programs and
expected results. Factored identities check the full construction;
independent derivations from all $19\,999$ nonzero ideal seeds end
at the nine additive generators. The archive retains deterministic
seed receipts, sample derivations, deliberately broken controls and
ordinary/optimized execution records. These supplement the universal
proofs; no enumeration of all $20\,000^3$ triples is asserted.
The universal nonpresentation lemma is proved above, rather than
inferred from finite examples.

AI tools were used extensively in construction, proof development,
computation, adversarial verification, literature research and writing.
This version is \textbf{unrefereed}. AI adversarial review is not external
human peer review, journal acceptance or formal proof-assistant
certification. The note and authored supporting materials are released
under CC BY 4.0.

\begin{thebibliography}{12}\small\raggedright
\bibitem{Bachiller}
D.~Bachiller,
Extensions, matched products, and simple braces,
\emph{J. Pure Appl. Algebra} \textbf{222} (2018), 1670--1691.
\href{https://arxiv.org/abs/1511.08477v3}{arXiv:1511.08477v3}.
\bibitem{OWR}
F.~Ced\'o,
Finite simple left braces: constructions and problems,
in \emph{Mini-Workshop: Algebraic Tools for Solving the Yang--Baxter Equation},
Oberwolfach Reports \textbf{51/2019}, 3226--3228.
\href{https://doi.org/10.4171/OWR/2019/51}{doi:10.4171/OWR/2019/51}.
\bibitem{CCS}
F.~Catino, I.~Colazzo and P.~Stefanelli,
Regular subgroups of the affine group and asymmetric product of radical braces,
\emph{J. Algebra} \textbf{455} (2016), 164--182.
\href{https://doi.org/10.1016/j.jalgebra.2016.01.038}{doi:10.1016/j.jalgebra.2016.01.038}.
\bibitem{BCJO}
D.~Bachiller, F.~Ced\'o, E.~Jespers and J.~Okni\'nski,
Asymmetric product of left braces and simplicity; new solutions of the
Yang--Baxter equation,
\emph{Commun. Contemp. Math.} \textbf{21} (2019), 1850042.
\href{https://arxiv.org/abs/1705.08493v1}{arXiv:1705.08493v1}.
\bibitem{CJOabundance}
F.~Ced\'o, E.~Jespers and J.~Okni\'nski,
An abundance of simple left braces with abelian multiplicative Sylow subgroups,
\emph{Rev. Mat. Iberoam.} \textbf{36} (2020), 1309--1332.
\href{https://doi.org/10.4171/rmi/1168}{doi:10.4171/rmi/1168}.
\bibitem{CJOembed}
F.~Ced\'o, E.~Jespers and J.~Okni\'nski,
Every finite abelian group is a subgroup of the additive group of a finite
simple left brace,
\emph{J. Pure Appl. Algebra} \textbf{225} (2021), 106476.
\href{https://arxiv.org/abs/2001.08905v1}{arXiv:2001.08905v1}.
\bibitem{Verbeken}
B.~Verbeken,
Powerful multiplicative groups do not force right nilpotence in finite braces,
\href{https://arxiv.org/abs/2608.01884v2}{arXiv:2608.01884v2} (2026).
\bibitem{CO}
F.~Ced\'o and J.~Okni\'nski,
New simple solutions of the Yang--Baxter equation and their permutation groups,
\emph{J. Pure Appl. Algebra} \textbf{229} (2025), 108025.
Inspected version:
\href{https://arxiv.org/abs/2401.12904v2}{arXiv:2401.12904v2}.
\bibitem{DameleKT}
M.~Damele,
Kernel--wreath constructions and infinite families of finite simple skew braces,
\href{https://arxiv.org/abs/2609.20401v1}{arXiv:2609.20401v1} (2026).
\bibitem{DameleIB}
M.~Damele,
Inverse Baer deformations and finite simple skew braces of nilpotent type,
\href{https://arxiv.org/abs/2609.24640v1}{arXiv:2609.24640v1} (2026).
\bibitem{Rump2020}
W.~Rump,
Affine structures of decomposable solvable groups,
\emph{J. Algebra} \textbf{556} (2020), 725--749.
\href{https://doi.org/10.1016/j.jalgebra.2020.04.004}{doi:10.1016/j.jalgebra.2020.04.004}.
Full body not inspected.
\bibitem{Rump2022}
W.~Rump,
Almost simple braces and primitive solutions to the Yang--Baxter equation,
\emph{J. Algebra} \textbf{601} (2022), 72--86.
\href{https://doi.org/10.1016/j.jalgebra.2022.03.008}{doi:10.1016/j.jalgebra.2022.03.008}.
Full body not inspected.
\end{thebibliography}
\end{document}
